\documentclass[11pt]{article}
\linespread{1.5}
\usepackage{fullpage}
\usepackage{amsmath,theorem,amssymb,graphicx,pgfplots,tabularx,placeins}
\usepackage[semicolon,authoryear]{natbib}
\usepackage{caption}
\usepackage{subcaption}
\usepackage{csquotes}
\usepackage{epstopdf}
\pgfplotsset{compat=1.18}

\newcommand{\lb}{\label}
\newtheorem{thm}{Theorem}
\newtheorem{prop}{Proposition}
\newtheorem{definition}{Definition}

\newcommand{\bit}{\begin{itemize}}
	\newcommand{\eit}{\end{itemize}}
\newcommand{\ben}{\begin{enumerate}}
	\newcommand{\een}{\end{enumerate}}
\newcommand\setItemnumber[1]{\setcounter{enumi}{\numexpr#1-1\relax}}

\title{Exam Ph.D. Macroeconomics II \\
Department of Economics, Uppsala University\\
June 10, 2026}
\author{}
\date{}

\begin{document}
\maketitle

\section*{Instructions}
\bit
	\item Writing time: 5 hours.
	
	\item The exam is closed book. 
	
	\item The exam has 70 points in total 
	
	\item A passing grade requires a) at least 30 points on  the exam, and b) 50 points in total for the course (incl the points you have from your problem sets; the problem set points will be rescaled so that they sum to 30 points: part I 25\%, part II 75\% of problem set points).
	
	\item Start each question on a new paper. Write your anonymous code on all answer pages.
	
	\item You may write your solutions by pen or pencil; use your best handwriting.
	
	\item Answers shall be given in English.

	\item Motivate your answers carefully; if you think you need to make additional assumptions to answer the questions, state them.
	
	\item If you have any questions during the exam, you may call Erik Öberg (+46 730 606 796) or Teodora Borota (+46 739 262 330) at any time between 3 PM and 5 PM.
\eit

\newpage

\section{Question 1: Optimal Subsidies in the Quality Ladders Model (20 points)}

Take the quality ladders model as studied in lectures, with some modifications. The physical marginal cost of making a machine of variety $\nu$ at quality $q$ is given by $\psi q(\nu,t \mid q)$. However, the government pays a subsidy $S_M \in (0,1)$ such that the actual marginal cost faced by machine producers is $\psi_e q(\nu,t \mid q)$, where $\psi_e \equiv (1-S_M)\psi.$

The government also pays a subsidy $S_{R\&D} \in (0,1)$ for R\&D expenditure. (Assume that the government budget is balanced by imposing lump-sum taxes on consumers.) Let us also set $\theta = 1$.\\

1. The final-good production function is $Y(t)=\frac{1}{1-\beta}L^\beta \int_0^1 q(\nu,t)x(\nu,t)^{1-\beta}\, d\nu .$ Show that the demand curve for a given variety is $x(\nu,t \mid q)=\left(\frac{q(\nu,t)}{p(\nu,t \mid q)}\right)^{1/\beta}L.$\\

2. The owner of the patent for variety $\nu$ at quality $q$ maximizes profits $\pi(\nu,t \mid q)=p(\nu,t \mid q)x(\nu,t \mid q)-\psi_e q(\nu,t \mid q)x(\nu,t \mid q).$ Imposing the normalization $1-\beta \equiv \psi$ from now on, show that $p(\nu,t \mid q)=(1-S_M)q(\nu,t)$ and $\pi(\nu,t \mid q)=(1-S_M)^{\frac{\beta-1}{\beta}}\beta q(\nu,t)L.$\\


3.Let $V(\nu,t \mid q)$ be the value of a patent that a successful innovator obtains if she raises the quality of variety $\nu$ to the level $q(\nu,t)$. Spending one unit of the final good on research generates a flow rate of innovation $\eta/q^I$, where $q^I=\frac{q(\nu,t)}{\lambda}$ is the incumbent's quality level. Interpret the condition for free entry into research, $1-S_{R\&D}=\frac{\eta V(\nu,t \mid q)}{q(\nu,t)/\lambda}.$\\

4. Interpret the no-arbitrage condition $r(t)V(\nu,t \mid q)=\pi(\nu,t \mid q)-z(\nu,t \mid q)V(\nu,t \mid q)+\dot V(\nu,t \mid q).$\\

5. Explain why under free entry, $\dot V(\nu,t \mid q)=0.$ Use the free entry and the no-arbitrage conditions, and the expression for profits you derived to show that $r(t)+z(\nu,t \mid q)=s\lambda\eta\beta L,$ where $s\equiv\frac{(1-S_M)^{\frac{\beta-1}{\beta}}}{1-S_{R\&D}}.$\\

6. On a balanced growth path (BGP), the following equations must hold: $r+z=s\lambda\eta\beta L,$ $g=r-\rho,$ and $g=(\lambda-1)z.$ Show that $g=\frac{s\lambda\eta\beta L-\rho}{(\lambda-1)^{-1}+1}.$\\

7. The socially optimal growth rate is $g^*=(\lambda-1)(1-\beta)^{-1/\beta}\eta\beta L-\rho.$ Show that to achieve this growth rate, it must be that $s=(1-\beta)^{-1/\beta}-\frac{\rho}{(\lambda-1)\lambda\eta\beta L}.$\\

8. True or false? Explain.
\begin{quote}
``Any combination of $S_M$ and $S_{R\&D}$ that satisfies equation $s=(1-\beta)^{-1/\beta}-\frac{\rho}{(\lambda-1)\lambda\eta\beta L}$ will achieve the socially optimal allocation.''
\end{quote}

\section{Government spending in the New-Keynesian model (20 points)}

Consider the vanilla New-Keynesian model studied in class, but now suppose that the government purchases a wasteful good $G_t$, financed by lump-sum taxes. Government purchases do not enter utility or production. The steady-state government share is $s_G \equiv \bar G/\bar Y \in (0,1)$.

The log-linearized household and firm optimality conditions are
\begin{eqnarray}
\hat c_t &=& -(\hat i_t-E_t\pi_{t+1})+E_t\hat c_{t+1}, \nonumber \\
\pi_t &=& \beta E_t\pi_{t+1}+\kappa \widehat{mc}_t, \nonumber \\
\hat \omega_t &=& \hat c_t+\varphi \hat n_t, \nonumber \\
\widehat{mc}_t &=& \hat\omega_t, \nonumber \\
\hat y_t &=& \hat n_t. \nonumber
\end{eqnarray}
Goods-market clearing is
\begin{eqnarray}
\hat y_t=(1-s_G)\hat c_t+s_G\hat g_t, \nonumber
\end{eqnarray}
where $\hat g_t$ follows
\begin{eqnarray}
\hat g_t=\rho_g \hat g_{t-1}+\epsilon^g_t, \qquad 0<\rho_g<1. \nonumber
\end{eqnarray}
Monetary policy is
\begin{eqnarray}
\hat i_t=\phi_\pi \pi_t, \qquad \phi_\pi>1. \nonumber
\end{eqnarray}

\ben
\item Derive an expression for real marginal cost as a function of $\hat y_t$ and $\hat g_t$. Then derive the flexible-price, or natural, level of output $\hat y^n_t$.

\item Define the output gap $x_t\equiv \hat y_t-\hat y^n_t$. Derive the two-equation New-Keynesian system in $(x_t,\pi_t)$ and derive the natural real interest rate $r^n_t$.

\item Guess that $x_t=a_x\hat g_t$ and $\pi_t=a_\pi\hat g_t$. Solve for $a_x$ and $a_\pi$. What are the signs of these coefficients? How does a larger $\phi_\pi$ affect the response of the output gap?

\item Compare the allocation under flexible prices, sticky prices with the Taylor rule above, and optimal monetary policy. In particular, explain whether output, private consumption, inflation, and the output gap increase or decrease after a positive government-spending shock.
\een


\section{Recruiting intensity in the DMP model (15 points)}

Consider the continuous-time DMP model studied in class, but suppose that firms posting vacancies can choose how intensively to recruit. A vacancy with recruiting intensity $e\geq 0$ meets workers at rate $e\lambda_v(\Theta)$, where $\Theta\equiv X/U$ is aggregate recruiting tightness and $X=eV$ is aggregate effective recruiting effort. Workers find jobs at rate $\lambda_u(\Theta)=\Theta\lambda_v(\Theta)$. Jobs separate at exogenous rate $\sigma$, employed workers produce $y$, unemployed workers receive flow utility $b$, and all agents discount at rate $r$. Entry is free.

Posting a vacancy has the standard flow cost $c$. In addition, recruiting intensity is costly. A firm choosing intensity $e$ pays the additional flow cost
\begin{eqnarray}
\frac{\kappa e^{1+\psi}}{1+\psi}, \qquad \kappa>0,\quad \psi>0. \nonumber
\end{eqnarray}
Wages are set by Nash bargaining with worker bargaining power $\gamma$. Individual firms take aggregate recruiting tightness $\Theta$ as given.

\ben
\item Write the Bellman equations for $W,U,J$ and $V$.

\item Define an equilibrium.

\item Derive the firm's optimal recruiting-intensity condition. How does the recruiting-intensity choice affect job-creation curve?

\item Derive the Nash-bargained wage equation. How does recruiting intensity affect the wage?

\item Suppose a new hiring technology lowers $\kappa$. What happens to recruiting intensity, recruiting tightness, unemployment, and wages? Explain the economics.
\een


\section{Idiosyncratic risk in the Aiyagari model (15 points)}

Consider a stationary Aiyagari economy with a continuum of ex-ante identical households. Each household has idiosyncratic labor efficiency $\epsilon_t$, drawn from a finite-state Markov process with invariant mean one. Households can save in a risk-free asset $a'$, but cannot borrow:
\begin{eqnarray}
a'\geq 0.
\end{eqnarray}
Preferences are
\begin{eqnarray}
E_0\sum_{t=0}^{\infty}\beta^t u(c_t),
\end{eqnarray}
where $u$ satisfies standard assumptions, including prudence. Firms are competitive and produce with
\begin{eqnarray}
Y=K^\alpha L^{1-\alpha},
\end{eqnarray}
with depreciation rate $\delta$.

\ben
\item Write the household's and firms' problem and derive their respective first-order conditions.

\item Define a stationary competitive equilibrium.

\item Draw and explain the Aiyagari diagram, with aggregate asset supply $A^s(r)$ and capital demand $K^d(r)$. Why must the equilibrium interest rate satisfy $r<1/\beta-1$?

\item Suppose idiosyncratic labor-income risk increases in the sense of a mean-preserving spread, while mean labor efficiency remains one. What happens to aggregate asset supply, the equilibrium interest rate, capital, and wages? Contrast this with a representative-agent or complete-markets economy.
\een

\end{document}
